\chapter{Extended Trajectory Derivations}\label{app:supp_trajectory}

This chapter provides the formal foundations and proofs for Chapter~\ref{ch:trajectory}.

\section{Layered Markov Process Foundations}\label{ptnb:sec:lmp-foundations}\label{subsec:trajectory_measures}

We collect here the layered-Markov-process abstraction that underlies every Hellinger statement in
this section: a finite layered graph augmented with a cemetery $\dagger$, and the act-state
\acp{ptnbMP} that the two games of this chapter instantiate on it. The natural \ac{rg} carries an
extra linear-routing state per neuron between consecutive activation states. We retain its
gate, branch, and routing substeps for full-law diagnostics; only ordinary surviving paths
admit exact act-only elimination, as detailed in Remark~\ref{ptnb:rem:rg-natural-equivalence}.
For full routing-process dynamic programming below, the layer index includes these refined
substeps; activation-boundary readouts form a subset of those layers.

\subsection{Layered Markov Process with Cemetery}\label{ptnb:sec:lmp}

\paragraph{State graph.}
Fix an integer $L \geq 1$.  For each $l \in \{0, 1, \ldots, L\}$ let $\mathcal{S}^{(l)}$ be a
finite set of \emph{ordinary states} at layer~$l$.  Adjoin a single absorbing \emph{cemetery state}
$\dagger$ shared across all layers and write $\mathcal{S}^{(l)}_+ \coloneqq \mathcal{S}^{(l)} \cup
    \{\dagger\}$.
The full state graph is the disjoint union $\bigsqcup_l \mathcal{S}^{(l)}_+$, and transitions are
only permitted from layer~$l$ to layer~$l{-}1$ (or into the cemetery, which then absorbs the
trajectory forever).  The states at layer~$0$, the input layer, are referred to as \emph{pixels}
below; the per-pixel decomposition will sum exactly over $\mathcal{S}^{(0)}_+$.

\paragraph{Transition kernel.}
For each layer $l \geq 1$, a \emph{backward transition kernel} $T^{(l)} : \mathcal{S}^{(l)}_+ \to
    \Delta\bigl(\mathcal{S}^{(l-1)}_+\bigr)$ assigns to each state $s \in \mathcal{S}^{(l)}_+$ a
probability distribution $T^{(l)}(s \to \cdot)$ over $\mathcal{S}^{(l-1)}_+$, with the constraint
that the cemetery is absorbing:
\begin{align}\label{ptnb:eq:cemetery-absorbing}
    T^{(l)}(\dagger \to \dagger) = 1, \qquad T^{(l)}(\dagger \to s') = 0 \;\;\text{for}\;\; s' \in
    \mathcal{S}^{(l-1)}.
\end{align}
Ordinary states may transition either to the next layer's ordinary states or into the cemetery
(early termination).  We write $T^{(l)}(s \to \dagger)$ for the probability that an ordinary state
at layer~$l$ jumps to the cemetery rather than to a layer-$(l{-}1)$ ordinary state.

\paragraph{Initial distribution and trajectories.}
Fix a deterministic starting state $s_0 \in \mathcal{S}^{(L)}$.  A \emph{trajectory} (indexed by
layer, $\zeta_l \in \mathcal{S}^{(l)}_+$) is a sequence
\begin{align}
    \zeta = (\zeta_L, \zeta_{L-1}, \ldots, \zeta_0) \;\in\; \mathcal{S}^{(L)}_+ \times
    \mathcal{S}^{(L-1)}_+ \times \cdots \times \mathcal{S}^{(0)}_+,
\end{align}
and the kernel $T^{(\cdot)}$ together with $\delta_{s_0}$ at layer~$L$ induces a probability
distribution on trajectories via
\begin{align}\label{ptnb:eq:path-measure}
    \Pi(\zeta_L, \zeta_{L-1}, \ldots, \zeta_0) \;=\; \mathbf{1}\{\zeta_L = s_0\} \prod_{k=1}^{L}
    T^{(k)}(\zeta_k \to \zeta_{k-1}).
\end{align}
For any $l \in \{0,\ldots,L\}$ the prefix distribution $\Pi^{(l)}$ over $(\zeta_L, \ldots,
    \zeta_l)$ is the marginal of~$\Pi$ over the trailing coordinates.

\paragraph{Two parallel processes.}
We compare two such processes that share the state graph and the start state~$s_0$, but use
different kernels: $M = A$ with $T^{(l)}_A, \Pi_A^{(l)}$, and $M = B$ with $T^{(l)}_B,
    \Pi_B^{(l)}$.
The kernels may differ either because the underlying network weights differ (e.g.\ randomization),
the activations differ (e.g.\ different inputs), or both.

\subsection{How the Two Games Fit This Abstraction}\label{ptnb:sec:games-fit}

We instantiate the layered MP of Section~\ref{ptnb:sec:lmp} for both games in this subsection. State
labels $s^{(l,\mathrm{type})}_{i,p,\kappa}$ follow Appendices~\ref{ptnb:app:relu-net-game}
and~\ref{ptnb:app:routing-formal}. We evaluate the SPNE policies of Theorems~\ref{ptnb:thm:forward-stopping}
(\ac{sg}) and~\ref{ptnb:thm:routing-forward} (\ac{rg}) at $\tau = 1$ (so activation-state decisions are
deterministic, continue iff $z_j^{(l)} > 0$); for the \ac{rg} we further fix $\varepsilon = 0$.

\paragraph{\ac{sg} MP.}
Ordinary states at layers $l \in \{1,\ldots,L\}$ are activation states
$s^{(l,\mathrm{act})}_{i,p}$, and layer-$0$ ordinary states are pixel-player terminal pairs
$s^{(0)}_{j,q}$ (each pixel splits into two terminals, mirroring the player tag the trajectory
carries to it). The deterministic start state is $s_0 = s^{(L,\mathrm{act})}_{u,+}$. The SPNE
backward kernel keeps the player on positive weights and flips it on negative weights:
\begin{align}
    T^{(l)}_M\bigl(s^{(l,\mathrm{act})}_{j,p} \to s^{(l-1,\mathrm{act})}_{i,p}\bigr)
    &= \frac{W^{(l,+)}_M[i,j]}{\gamma^{(l)}_{M,j}}\,\1\{z_j^{(l)} > 0\},
    \nonumber \\
    T^{(l)}_M\bigl(s^{(l,\mathrm{act})}_{j,p} \to s^{(l-1,\mathrm{act})}_{i,p'}\bigr)
    &= \frac{W^{(l,-)}_M[i,j]}{\gamma^{(l)}_{M,j}}\,\1\{z_j^{(l)} > 0\},
    \qquad 2 \le l \le L,
    \label{ptnb:eq:adf-T}
\end{align}
with $\gamma^{(l)}_{M,j} = \sum_i |W^{(l)}_M[i,j]|$, $p' = -p$, cemetery row
$T^{(l)}_M(s^{(l,\mathrm{act})}_{j,p}\to\dagger)=1-\1\{z_j^{(l)}>0\}\1\{\gamma^{(l)}_{M,j}>0\}$, with $0/0\coloneqq0$ in the ordinary transitions, and the analogous form
at $l = 1$ landing on the pixel-player terminals $s^{(0)}_{i,p}, s^{(0)}_{i,p'}$. This is the
natural \ac{sg} MP of Section~\ref{ptnb:app:relu-net-game} verbatim.

\paragraph{\ac{rg} MP.}
Ordinary states at layers $l \in \{1,\ldots,L\}$ are activation states
$s^{(l,\mathrm{act})}_{i,p}$, and layer-$0$ ordinary states are pixel-stream terminal pairs
$s^{(0)}_{j,\kappa}$ (each pixel splits into two terminals tagged by the layer-$1$ routing-edge
stream). The deterministic start state is $s_0 = s^{(L,\mathrm{act})}_{u,+}$. The kernel composes
the activation-state stop/continue decision and the subsequent $\kappa$-branched Gibbs routing into
one act-to-act transition:
\begin{align}
    T^{(l)}_M\bigl(s^{(l,\mathrm{act})}_{j,p} \to s^{(l-1,\mathrm{act})}_{i,p}\bigr)
    &= \tfrac12\,\pi^{(l)}_{M,+}(i \mid j)\,\1\{z_j^{(l)} > 0\},
    \nonumber \\
    T^{(l)}_M\bigl(s^{(l,\mathrm{act})}_{j,p} \to s^{(l-1,\mathrm{act})}_{i,p'}\bigr)
    &= \tfrac12\,\pi^{(l)}_{M,-}(i \mid j)\,\1\{z_j^{(l)} > 0\},
    \qquad 2 \le l \le L,
    \label{ptnb:eq:rg-T-act}
\end{align}
and at $l = 1$ the destinations are pixel-stream terminals carrying the routing-edge $\kappa$,
\begin{align}\label{ptnb:eq:rg-T-pixel}
    T^{(1)}_M\bigl(s^{(1,\mathrm{act})}_{j,p} \to s^{(0)}_{i,\kappa}\bigr)
    \;=\; \tfrac12\,\pi^{(1)}_{M,\kappa}(i \mid j)\,\1\{z_j^{(1)} > 0\},
    \qquad \kappa \in \{+,-\},
\end{align}
For this act-only marginal, the cemetery probability is the residual
\[
T^{(l)}_M(s_{j,p}\to\dagger)=1-\sum_{s'\in\mathcal S^{(l-1)}}T^{(l)}_M(s_{j,p}\to s').
\]
An empty routing stream has zero ordinary policy mass and its $1/2$ branch share
is absorbed, even when the source neuron is alive. The Gibbs policy on a nonempty
stream is the temperature-$\tau$
distribution
\begin{align}\label{ptnb:eq:rg-T}
    \pi^{(l)}_{M,\kappa}(i \mid j)
    \;\propto\;
    \bigl[\,a^{(l-1)}_{M,i}\;W^{(l,\kappa)}_{M,ij}\,\bigr]^{1/\tau},
    \qquad i \in \mathrm{Pred}_j^{(l)},
\end{align}
(at $l = 1$, $a_{M,i}^{(0)}\,W^{(1,\kappa)}_{M,ij}$ is replaced by $[x_i\,W^{(1)}_{M,ij}]^\kappa$;
when $\varepsilon>0$, each continuing $\kappa$ branch sends its
$\tfrac12\varepsilon^{1/\tau}/Z^{(l)}_{j,\kappa}$ share to the cemetery; as in the
implementation, the bias node of Section~\ref{ptnb:app:routing-formal} is omitted).
If all scores in a stream are $-\infty$, its ordinary policy is zero and the whole
branch is absorbed, in agreement with the outside-option convention.
Equations~\eqref{ptnb:eq:rg-T-act}--\eqref{ptnb:eq:rg-T-pixel} are the act-to-act marginal of the \emph{diagnostic} routing process: it uses the original activation gates, omits the bias action, and normalizes each remaining routing branch. When the original bias mass is nonzero, this differs from the bias-inclusive game of Section~\ref{ptnb:app:routing-formal}. For example, $a=W=b=1$, $\varepsilon=0$, gives probability $1/4$ on the ordinary positive route in that game, but $1/2$ after omitting the bias action. Remark~\ref{ptnb:rem:rg-natural-equivalence} concerns state elimination within the same chosen process; it does not identify these two different kernels. The reported full-trajectory diagnostics use
the refined gate, branch, and routing substeps: absorption at different substeps
is retained in the trajectory history. The ordinary surviving-path overlap can
be computed from the displayed act-to-act kernels, but their collapsed cemetery
row must not be substituted for the refined process when computing the full overlap.

\subsection{Marginal Occupation Measures}\label{ptnb:sec:lmp-marginals}

For each model $M \in \{A,B\}$, the marginal probability that the chain visits state $s$ at
layer~$l$ is
\begin{align}\label{ptnb:eq:lmp-alpha}
    m^{(l)}_M(s) \;=\; \sum_{\zeta: \zeta_l = s} \Pi^{(l)}_M(\zeta),
    \qquad s \in \mathcal{S}^{(l)}_+.
\end{align}
This satisfies the deterministic boundary condition
\begin{align}\label{ptnb:eq:lmp-alpha-boundary}
    m^{(L)}_M(s_0) = 1, \qquad m^{(L)}_M(s) = 0 \;\;\text{for}\;\; s \neq s_0,
\end{align}
and the one-step recursion
\begin{align}\label{ptnb:eq:lmp-alpha-recursion}
    m^{(l-1)}_M(s') \;=\; \sum_{s \in \mathcal{S}^{(l)}_+} m^{(l)}_M(s)\; T^{(l)}_M(s \to s'),
    \qquad s' \in \mathcal{S}^{(l-1)}_+.
\end{align}
We refer to the single backward propagation~\eqref{ptnb:eq:lmp-alpha-recursion} as the \emph{Marginal
Pass}.  At the input layer, $m^{(0)}_M(s)$ is the probability that a backward trajectory from~$s_0$
terminates at pixel~$s$ under model~$M$, and $m^{(0)}_M(\dagger)$ is the total mass absorbed by the
cemetery along the way.

\section{Hellinger Distance from Bhattacharyya Coefficients}\label{ptnb:sec:lmp-hellinger-distance}\label{subsec:hellinger_definition}

\paragraph{Properties.}
$\mathrm{BC}^{(l)}$ and $\mathrm{H}^{(l)}$ lie in $[0,1]$ for every $l$, even under support
mismatch.  If a state is alive in one model and dead in the other, its geometric-mean row vanishes
and its mass contributes directly to the Hellinger increment.  Since $\mathrm{H}$ is a metric and
satisfies $\mathrm{H}^2 \le d_{\mathrm{TV}} \le \sqrt{2}\,\mathrm{H}$ and $\mathrm{H}^2 \le
    \tfrac{1}{2}\mathrm{KL}$, small Hellinger distance controls TV and small KL controls H as well.
\emph{Normalization conventions.} Here $d_{\mathrm{TV}}(\Pi_A, \Pi_B) = \tfrac{1}{2}\sum_{\zeta}
    |\Pi_A(\zeta)-\Pi_B(\zeta)| \in [0,1]$ is the probability-metric (half-$L^1$) total-variation
distance, and $\mathrm{KL}$ is measured in nats; the displayed inequalities follow this convention.
Readers using the unnormalized $L^1$ total variation $\|\Pi_A-\Pi_B\|_1 = 2 d_{\mathrm{TV}}$ should
absorb the factor of~$2$ accordingly.

\paragraph{Alternative divergences.}\phantomsection\label{ptnb:app:why-hellinger}
KL and Jeffreys divergence have an exact Markov-chain chain rule but are unbounded and become
infinite under support mismatch, e.g.\ when randomization opens or closes routes.  Jensen--Shannon
is bounded and finite, but $\tfrac{1}{2}(\Pi_A + \Pi_B)$ is not Markov on the layered graph, so no
exact linear-time dynamic program is available.  Wasserstein distances need a ground cost between
neuron-labeled trajectories (across layers, signs and permuted neurons), for which there is no
canonical choice, and the resulting transport problem has no exact local recursion.  Hellinger is
bounded, finite, and exactly tractable: the overlap factorizes along the layered law
(Lemma~\ref{ptnb:lem:bc-chain-rule}), which gives a dynamic program linear in the number of edges
(Algorithm~\ref{ptnb:alg:hellinger}) and the non-negative per-state decomposition of
Theorem~\ref{ptnb:thm:per-pixel}.  Note that $\mathrm{H}$ compares the \emph{undiscounted} law $\Pi$:
the discount $d_t$ does not enter the kernels, so $\mathrm{H}$ is blind to the path weights
$(\alpha,\beta)$ and to the \ac{sg} weight magnitudes $\gamma^{(l)}_j$, which enter only through
$d_t$; it compares \emph{where} trajectories go, not how much weight they carry.

\paragraph{Computation.}
Compute the marginal occupations $m_M^{(l)}$, run the Hellinger pass
(Algorithm~\ref{ptnb:alg:hellinger}), and read off $\mathrm{H}^2(s)$ from~\eqref{ptnb:eq:per-pixel-h2}; the
total cost is linear in the number of edges.

\section{Hellinger Identities and Proofs}

\subsection{Multiplicative Chain Rule}\label{ptnb:sec:lmp-chain-rule}

\begin{lemmaenv}[Chain rule, \LeanResult{ptnb:lem:bc-chain-rule}]\label{ptnb:lem:bc-chain-rule}
For every $l \in \{0, 1, \ldots, L\}$, the Bhattacharyya coefficient of the prefix distributions
factorizes along trajectories:
\begin{align}\label{ptnb:eq:bc-chain-rule}
    \mathrm{BC}\!\bigl(\Pi^{(l)}_A,\, \Pi^{(l)}_B\bigr) \;=\; \sum_{(\zeta_L, \ldots, \zeta_l)}
    \mathbf{1}\{\zeta_L = s_0\} \prod_{k=l+1}^{L} \widetilde{T}^{(k)}(\zeta_k, \zeta_{k-1}).
\end{align}\end{lemmaenv}

\begin{proof}
Substitute the Markov factorization $\Pi^{(l)}_M(\zeta) = \mathbf{1}\{\zeta_L =
    s_0\}\prod_{k=l+1}^L T^{(k)}_M(\zeta_k, \zeta_{k-1})$ for $M \in \{A,B\}$ into the definition of
$\mathrm{BC}$:
\begin{align}
    \mathrm{BC}\!\bigl(\Pi^{(l)}_A,\, \Pi^{(l)}_B\bigr)
    &\;=\; \sum_{(\zeta_L, \ldots, \zeta_l)} \mathbf{1}\{\zeta_L = s_0\}\,
    \sqrt{\prod_{k=l+1}^{L} T^{(k)}_A(\zeta_k, \zeta_{k-1})\, T^{(k)}_B(\zeta_k, \zeta_{k-1})} \\
    &\;=\; \sum_{(\zeta_L, \ldots, \zeta_l)} \mathbf{1}\{\zeta_L = s_0\}\,
    \prod_{k=l+1}^{L} \widetilde{T}^{(k)}(\zeta_k, \zeta_{k-1}),
\end{align}
using $\sqrt{\prod_k a_k b_k} = \prod_k \sqrt{a_k b_k}$ for non-negative factors.
\end{proof}

\begin{lemmaenv}[Monotonicity, \LeanResult{ptnb:lem:bc-monotone}]\label{ptnb:lem:bc-monotone}
Let $\mathrm{BC}^{(l)} \coloneqq \mathrm{BC}(\Pi^{(l)}_A, \Pi^{(l)}_B)$.  Then $\mathrm{BC}^{(l)}$
is monotonically non-increasing as $l$ decreases from $L$ to $0$:
\begin{align}
    1 \;=\; \mathrm{BC}^{(L)} \;\geq\; \mathrm{BC}^{(L-1)} \;\geq\; \cdots \;\geq\; \mathrm{BC}^{(0)}
    \;\geq\; 0.
\end{align}
Equivalently, the squared Hellinger distance $\mathrm{H}^{2}_{(l)} = 1 - \mathrm{BC}^{(l)}$ is
monotonically non-decreasing.\end{lemmaenv}

\begin{proof}
Define
\begin{align}
    g^{(l)}(s) \;\coloneqq\; \sum_{(\zeta_L,\ldots,\zeta_l)} \mathbf{1}\{\zeta_L = s_0\}\,
    \prod_{k=l+1}^L \widetilde{T}^{(k)}(\zeta_k, \zeta_{k-1})\, \mathbf{1}\{\zeta_l = s\}.
\end{align}
From Lemma~\ref{ptnb:lem:bc-chain-rule}, $\mathrm{BC}^{(l)} = \sum_{s \in \mathcal{S}^{(l)}_+}
    g^{(l)}(s)$.  Grouping the layer-$(l-1)$ sum by $\zeta_l = s$,
\begin{align}
    \mathrm{BC}^{(l-1)} \;=\; \sum_{s \in \mathcal{S}^{(l)}_+} g^{(l)}(s)\; \widetilde{R}^{(l)}(s).
\end{align}
Since $\widetilde{R}^{(l)}(s) \in [0, 1]$ and $g^{(l)}(s) \geq 0$, we get $\mathrm{BC}^{(l-1)} \leq
    \mathrm{BC}^{(l)}$.  The base case $\mathrm{BC}^{(L)} = 1$ holds because $\Pi^{(L)}_A = \Pi^{(L)}_B
    = \delta_{s_0}$ and $\sqrt{1 \cdot 1} = 1$.
\end{proof}

\begin{lemmaenv}[Backward pass correctness, \LeanResult{ptnb:lem:hellinger-backward-correctness}]\label{ptnb:lem:hellinger-backward-correctness}
Algorithm~\ref{ptnb:alg:hellinger} computes $g^{(l)}$, hence $\mathrm{BC}^{(l)}$, $\mathrm{H}^{(l)}$,
and $\mathrm{H}^2(s)$ correctly for every $l$ and $s$.  In particular,
\begin{align}
    \sum_{s \in \mathcal{S}^{(0)}_+} \mathrm{H}^2(s) \;=\; \mathrm{H}^{2}_{(0)}.
\end{align}\end{lemmaenv}

\begin{proof}
\emph{Correctness of $g^{(l)}$ by downward induction on $l$.}  At $l = L$, the initialization
$g^{(L)}(s_0) = 1$ and $g^{(L)}(s) = 0$ for $s \neq s_0$ matches the definition $g^{(L)}(s) =
    \mathbf{1}\{s = s_0\}$ in Lemma~\ref{ptnb:lem:bc-monotone} (the empty product equals~1).

Assume $g^{(l)}$ stored by the algorithm equals the $g^{(l)}$ defined in
Lemma~\ref{ptnb:lem:bc-monotone}.  In iteration $l$, the algorithm computes
\begin{align}
    g^{(l-1)}(s') \;=\; \sum_{s \in \mathcal{S}^{(l)}_+} g^{(l)}(s)\, \widetilde{T}^{(l)}(s, s'),
\end{align}
which by direct expansion equals
\begin{align}
    \begin{aligned}
        &\sum_s \Bigl(\sum_{(\zeta_L,\ldots,\zeta_l)} \mathbf{1}\{\zeta_L =
        s_0\}\!\!\prod_{k=l+1}^L\!\!\widetilde{T}^{(k)}(\zeta_k, \zeta_{k-1})\, \mathbf{1}\{\zeta_l =
        s\}\Bigr) \widetilde{T}^{(l)}(s, s') \\
        &\qquad=\; \sum_{(\zeta_L, \ldots, \zeta_{l-1})} \mathbf{1}\{\zeta_L =
        s_0\}\!\!\prod_{k=l}^L\!\!\widetilde{T}^{(k)}(\zeta_k, \zeta_{k-1})\, \mathbf{1}\{\zeta_{l-1} =
        s'\}.
    \end{aligned}
\end{align}
This equals $g^{(l-1)}(s')$ from Lemma~\ref{ptnb:lem:bc-monotone}, closing the induction.

\emph{Correctness of $\mathrm{BC}^{(l)}$.}  After the loop, the algorithm sets $\mathrm{BC}^{(l)} =
    \sum_{s'} g^{(l)}(s')$, which by the proof of Lemma~\ref{ptnb:lem:bc-monotone} equals
$\mathrm{BC}\!\bigl(\Pi^{(l)}_A, \Pi^{(l)}_B\bigr)$.  Then $\mathrm{H}^{(l)} = \sqrt{1 -
    \mathrm{BC}^{(l)}}$ by definition~\eqref{ptnb:eq:hellinger-def}.

\emph{Correctness of $\mathrm{H}^2(\cdot)$.} At $l = 0$, the stored value satisfies
\begin{align}
    g^{(0)}(s)
    =
    \sum_{\zeta: \zeta_0 = s}
    \sqrt{\Pi^{(0)}_A(\zeta) \Pi^{(0)}_B(\zeta)},
\end{align}
which is exactly the quantity used in the proof of Theorem~\ref{ptnb:thm:per-pixel}. Substituting
into~\eqref{ptnb:eq:per-pixel-h2} gives the per-pixel disagreement values from
Theorem~\ref{ptnb:thm:per-pixel}, and~\eqref{ptnb:eq:per-pixel-sum} guarantees they sum to
$\mathrm{H}^{2}_{(0)}$.
\end{proof}

\begin{remark}[Routing-state elimination and absorption histories, \LeanResult{ptnb:rem:rg-natural-equivalence}]\label{ptnb:rem:rg-natural-equivalence}
The refined diagnostic routing process just defined alternates activation and linear-routing states.
A dead activation stops at the gate substep; an alive activation chooses its
$\kappa$ branch uniformly, after which an empty branch or outside option can
absorb at the routing substep. The refined layered process retains these
substeps and pads the cemetery state after absorption. Algorithm~\ref{ptnb:alg:hellinger}
applies to this expanded state graph; the diagnostics are read at the activation
boundaries. This is the process implemented by the separate gate and branch
absorption updates.

For a trajectory that remains ordinary, deleting each intermediate routing
state is a bijection: its neuron and player are determined by the adjacent
activation states, and its branch by whether the player is kept or switched
(or by the pixel-stream tag at the input). The act-to-act formulas therefore
preserve ordinary surviving-path probabilities and their Bhattacharyya overlap.
The same applies after conditioning both laws on survival.

For absorbing trajectories, deletion need not be injective: a gate stop and a
later empty-branch stop can become the same act-only trajectory. Thus full-law
Bhattacharyya and Hellinger equality is asserted only when absorption histories
are recoverable, for example when no alive routing branch leaks to the cemetery.
In general the finite-fibre Cauchy--Schwarz inequality gives
\[
 \mathrm{BC}_{\mathrm{refined}}\leq\mathrm{BC}_{\mathrm{act\text{-}only}},
 \qquad
 \mathrm{H}_{\mathrm{refined}}\geq\mathrm{H}_{\mathrm{act\text{-}only}}.
\]
For example, take one input $x=1$ and weight $W=1$. Model $A$ has bias $0$ and
continues, with its empty negative branch absorbing with probability $1/2$;
model $B$ has bias $-2$ and stops at the gate. Their refined path supports are
disjoint, so $\mathrm{BC}_{\mathrm{refined}}=0$. Collapsing the two absorption
substeps instead gives $\mathrm{BC}_{\mathrm{act\text{-}only}}=1/\sqrt2$.\end{remark}

\section{Proof of the Per-Pixel Decomposition}

The following proves Theorem~\ref{ptnb:thm:per-pixel}.

\begin{proof}
Apply the algebraic identity $(\sqrt{a} - \sqrt{b})^2 = a + b - 2\sqrt{ab}$ to each trajectory:
\begin{align}
    \mathrm{H}^2 \;=\; \tfrac{1}{2}\sum_{\zeta}\bigl(\sqrt{\Pi^{(0)}_A(\zeta)} -
    \sqrt{\Pi^{(0)}_B(\zeta)}\bigr)^2
    \;=\; \tfrac{1}{2}\sum_{\zeta}\bigl[\Pi^{(0)}_A(\zeta) + \Pi^{(0)}_B(\zeta) -
    2\sqrt{\Pi^{(0)}_A(\zeta)\Pi^{(0)}_B(\zeta)}\bigr].
\end{align}
Group the trajectory sum by the terminal pixel $\zeta_0 = s$:
\begin{align}
    \mathrm{H}^2 \;=\; \tfrac{1}{2}\sum_{s \in \mathcal{S}^{(0)}_+}\Bigl[\sum_{\zeta: \zeta_0 =
    s}\Pi^{(0)}_A(\zeta)
    + \sum_{\zeta: \zeta_0 = s}\Pi^{(0)}_B(\zeta)
    - 2\sum_{\zeta: \zeta_0 = s}\sqrt{\Pi^{(0)}_A(\zeta)\Pi^{(0)}_B(\zeta)}\Bigr].
\end{align}
By definition, $\sum_{\zeta: \zeta_0 = s}\Pi^{(0)}_M(\zeta) = m^{(0)}_M(s)$ for $M \in \{A,B\}$.
By Lemma~\ref{ptnb:lem:bc-chain-rule} restricted to trajectories terminating at~$s$ (i.e.\ tracing
through Lemma~\ref{ptnb:lem:bc-monotone}'s definition of $g^{(l)}$), $\sum_{\zeta: \zeta_0 =
    s}\sqrt{\Pi^{(0)}_A(\zeta)\Pi^{(0)}_B(\zeta)} = g^{(0)}(s)$.  Substituting,
\begin{align}
    \mathrm{H}^2 \;=\; \tfrac{1}{2}\sum_s\bigl[m^{(0)}_A(s) + m^{(0)}_B(s) - 2g^{(0)}(s)\bigr]
    \;=\; \sum_{s \in \mathcal{S}^{(0)}_+} \mathrm{H}^2(s),
\end{align}
proving (ii).

For (i), apply Cauchy--Schwarz to the sum defining $g^{(0)}(s)$:
\begin{align}
    g^{(0)}(s)
    = \sum_{\zeta: \zeta_0 = s} \sqrt{\Pi^{(0)}_A(\zeta) \Pi^{(0)}_B(\zeta)}
    \leq \sqrt{\Bigl(\sum_{\zeta: \zeta_0 = s}\Pi^{(0)}_A(\zeta)\Bigr)\Bigl(\sum_{\zeta: \zeta_0 =
    s}\Pi^{(0)}_B(\zeta)\Bigr)}
    = \sqrt{m^{(0)}_A(s)\,m^{(0)}_B(s)}.
\end{align}
Combined with the AM--GM inequality $\sqrt{m^{(0)}_A(s)m^{(0)}_B(s)} \leq \tfrac{1}{2}(m^{(0)}_A(s)
    + m^{(0)}_B(s))$, this gives $g^{(0)}(s) \leq \tfrac{1}{2}(m^{(0)}_A(s) + m^{(0)}_B(s))$, so
$\mathrm{H}^2(s) \geq 0$.  Equality holds iff $\Pi^{(0)}_A(\cdot \mid \zeta_0 = s) =
    \Pi^{(0)}_B(\cdot \mid \zeta_0 = s)$ \emph{and} $m^{(0)}_A(s) = m^{(0)}_B(s)$.

Item (iii) is immediate from (ii) by separating the cemetery from the ordinary pixels.
\end{proof}

\section{Conditioned-Survival Trajectory Laws}

\begin{theoremenv}[Conditioned-survival trajectory law and conditioned Hellinger distance, \LeanResult{ptnb:thm:conditioned-survival}]\label{ptnb:thm:conditioned-survival}
Let $\mathsf{Surv} \coloneqq \{(\zeta_L,\ldots,\zeta_0) : \zeta_0 \in \mathcal{S}^{(0)}\}$ be the
event, in the trajectory space $\prod_{l=0}^{L} \mathcal{S}^{(l)}_+$, that the backward trajectory
ends in an ordinary input-layer state — equivalently, the event that the trajectory never enters
the cemetery~$\dagger$.  For each model $M \in \{A,B\}$ define the survival function
\begin{align}\label{ptnb:eq:survival-function}
    \psi_M^{(l)}(s)
    \;\coloneqq\;
    \mathbb{P}_M(\mathsf{Surv} \mid \zeta_l = s),
    \qquad s \in \mathcal{S}^{(l)}_+, \;\; l \in \{0,\ldots,L\}.
\end{align}
Then:
\begin{enumerate}
    \item[(i)] The boundary and recursion are
        \begin{align}\label{ptnb:eq:survival-boundary}
            \psi_M^{(0)}(s) = 1 \;\; \text{for } s \in \mathcal{S}^{(0)},
            \qquad
            \psi_M^{(0)}(\dagger) = 0,
        \end{align}
        and for every $l \ge 1$,
        \begin{align}\label{ptnb:eq:survival-recursion}
            \psi_M^{(l)}(s)
            \;=\;
            \sum_{s' \in \mathcal{S}^{(l-1)}_+} T_M^{(l)}(s \to s')\, \psi_M^{(l-1)}(s')
            \;=\;
            \sum_{s' \in \mathcal{S}^{(l-1)}} T_M^{(l)}(s \to s')\, \psi_M^{(l-1)}(s').
        \end{align}
        In particular,
        \begin{align}\label{ptnb:eq:survival-mass}
            Z_M \;\coloneqq\; \mathbb{P}_M(\mathsf{Surv}) \;=\; \psi_M^{(L)}(s_0) \;=\; \sum_{s \in
            \mathcal{S}^{(0)}} m_M^{(0)}(s) \;=\; 1 - m_M^{(0)}(\dagger).
        \end{align}
    \item[(ii)] Assume $Z_M > 0$.  For every ordinary state $s \in \mathcal{S}^{(l)}$ with $\psi_M^{(l)}(s) >
            0$, define the conditioned kernel on ordinary states by
        \begin{align}\label{ptnb:eq:conditioned-kernel}
            T_{M,\mathrm{surv}}^{(l)}(s \to s')
            \;\coloneqq\;
            \frac{T_M^{(l)}(s \to s')\, \psi_M^{(l-1)}(s')}{\psi_M^{(l)}(s)},
            \qquad s' \in \mathcal{S}^{(l-1)}.
        \end{align}
        Then $T_{M,\mathrm{surv}}^{(l)}(s \to \cdot)$ is row-stochastic on $\mathcal{S}^{(l-1)}$, and the
        conditional trajectory distribution $\Pi_{M,\mathrm{surv}} \coloneqq \Pi_M(\,\cdot \mid
            \mathsf{Surv}\,)$ factorizes as the layered Markov law
        \begin{align}\label{ptnb:eq:conditioned-path-measure}
            \Pi_{M,\mathrm{surv}}(\zeta_L,\ldots,\zeta_0)
            = \mathbf{1}\{\zeta_L = s_0\} \prod_{k=1}^{L} T_{M,\mathrm{surv}}^{(k)}(\zeta_k \to \zeta_{k-1}),
            \qquad (\zeta_L,\ldots,\zeta_0) \in \mathsf{Surv}.
        \end{align}
        Its layer-$l$ marginals satisfy
        \begin{align}\label{ptnb:eq:conditioned-alpha}
            m_{M,\mathrm{surv}}^{(l)}(s)
            \;=\;
            \frac{m_M^{(l)}(s)\, \psi_M^{(l)}(s)}{Z_M},
            \qquad s \in \mathcal{S}^{(l)}.
        \end{align}
    \item[(iii)] Define the conditioned geometric-mean kernel
        \begin{align}\label{ptnb:eq:conditioned-geom-kernel}
            \widetilde{T}_{\mathrm{surv}}^{(l)}(s,s')
            \coloneqq \sqrt{T_{A,\mathrm{surv}}^{(l)}(s \to s')\, T_{B,\mathrm{surv}}^{(l)}(s \to s')}
            = \frac{\widetilde{T}^{(l)}(s,s')\, \sqrt{\psi_A^{(l-1)}(s')\,
            \psi_B^{(l-1)}(s')}}{\sqrt{\psi_A^{(l)}(s)\, \psi_B^{(l)}(s)}}.
        \end{align}
        The Bhattacharyya coefficient of the conditioned trajectory laws is
        \begin{align}\label{ptnb:eq:conditioned-bc-kernel}
            \mathrm{BC}\!\bigl(\Pi_{A,\mathrm{surv}}, \Pi_{B,\mathrm{surv}}\bigr)
            \;=\;
            \sum_{(\zeta_L,\ldots,\zeta_0) \in \mathsf{Surv}}
            \mathbf{1}\{\zeta_L = s_0\}
            \prod_{k=1}^{L} \widetilde{T}_{\mathrm{surv}}^{(k)}(\zeta_k, \zeta_{k-1}),
        \end{align}
        so it is computed by the same Hellinger pass (Algorithm~\ref{ptnb:alg:hellinger}), but on the
        ordinary-state graph with the conditioned kernel.
    \item[(iv)] Let
        \begin{align}\label{ptnb:eq:ordinary-bc}
            \mathrm{BC}_{\mathrm{ord}}^{(0)}
            \;\coloneqq\;
            \sum_{s \in \mathcal{S}^{(0)}} g^{(0)}(s)
            \;=\;
            \sum_{(\zeta_L,\ldots,\zeta_0) \in \mathsf{Surv}} \sqrt{\Pi_A(\zeta_L,\ldots,\zeta_0)\,
            \Pi_B(\zeta_L,\ldots,\zeta_0)}.
        \end{align}
        Then
        \begin{align}\label{ptnb:eq:conditioned-bc-posthoc}
            \mathrm{BC}\!\bigl(\Pi_{A,\mathrm{surv}}, \Pi_{B,\mathrm{surv}}\bigr)
            \;=\;
            \frac{\mathrm{BC}_{\mathrm{ord}}^{(0)}}{\sqrt{Z_A Z_B}},
        \end{align}
        and the conditioned-survival Hellinger distance is
        \begin{align}\label{ptnb:eq:conditioned-h2}
            \mathrm{H}_{\mathrm{surv}}^2
            \;\coloneqq\;
            \mathrm{H}^2\!\bigl(\Pi_{A,\mathrm{surv}}, \Pi_{B,\mathrm{surv}}\bigr)
            \;=\;
            1 - \frac{\mathrm{BC}_{\mathrm{ord}}^{(0)}}{\sqrt{Z_A Z_B}}.
        \end{align}
\end{enumerate}\end{theoremenv}

\begin{proof}
For (i),~\eqref{ptnb:eq:survival-boundary} is immediate from the definition of $\mathsf{Surv}$: at
layer~$0$ every ordinary state survives and the cemetery does not.  For $l \ge 1$, conditioning on
the next step gives
\begin{align}
    \psi_M^{(l)}(s)
    =
    \sum_{s' \in \mathcal{S}^{(l-1)}_+}
    T_M^{(l)}(s \to s')\, \mathbb{P}_M(\mathsf{Surv} \mid \zeta_{l-1} = s')
    =
    \sum_{s' \in \mathcal{S}^{(l-1)}_+}
    T_M^{(l)}(s \to s')\, \psi_M^{(l-1)}(s').
\end{align}
Since $\psi_M^{(l-1)}(\dagger)=0$, the cemetery term drops out and~\eqref{ptnb:eq:survival-recursion}
follows.  Evaluating at $s_0$ gives $Z_M = \psi_M^{(L)}(s_0)$.  The identities with $m_M^{(0)}$ are
just the layer-$0$ marginal of the event $\mathsf{Surv}$.

For (ii), row-stochasticity follows from~\eqref{ptnb:eq:survival-recursion}:
\begin{align}
    \sum_{s' \in \mathcal{S}^{(l-1)}} T_{M,\mathrm{surv}}^{(l)}(s \to s')
    =
    \frac{1}{\psi_M^{(l)}(s)}
    \sum_{s' \in \mathcal{S}^{(l-1)}} T_M^{(l)}(s \to s')\, \psi_M^{(l-1)}(s')
    = 1.
\end{align}
Now fix $(\zeta_L,\ldots,\zeta_0) \in \mathsf{Surv}$.  By definition of conditional probability,
\begin{align}
    \Pi_{M,\mathrm{surv}}(\zeta_L,\ldots,\zeta_0)
    =
    \frac{\Pi_M(\zeta_L,\ldots,\zeta_0)}{Z_M}
    =
    \frac{\mathbf{1}\{\zeta_L=s_0\}}{Z_M}
    \prod_{k=1}^{L} T_M^{(k)}(\zeta_k \to \zeta_{k-1}).
\end{align}
Multiplying the conditioned kernels and telescoping the survival factors gives
\begin{align}
    \prod_{k=1}^{L} T_{M,\mathrm{surv}}^{(k)}(\zeta_k \to \zeta_{k-1})
    &=
    \prod_{k=1}^{L}
    T_M^{(k)}(\zeta_k \to \zeta_{k-1})
    \frac{\psi_M^{(k-1)}(\zeta_{k-1})}{\psi_M^{(k)}(\zeta_k)} \\
    &=
    \frac{\psi_M^{(0)}(\zeta_0)}{\psi_M^{(L)}(s_0)}
    \prod_{k=1}^{L} T_M^{(k)}(\zeta_k \to \zeta_{k-1}) \\
    &=
    \frac{1}{Z_M}
    \prod_{k=1}^{L} T_M^{(k)}(\zeta_k \to \zeta_{k-1}),
\end{align}
because $\zeta_0 \in \mathcal{S}^{(0)}$ implies $\psi_M^{(0)}(\zeta_0)=1$.  This
is~\eqref{ptnb:eq:conditioned-path-measure}.  For the marginals,
\begin{align}
    m_{M,\mathrm{surv}}^{(l)}(s)
    =
    \mathbb{P}_M(\zeta_l = s \mid \mathsf{Surv})
    =
    \frac{\mathbb{P}_M(\zeta_l = s, \mathsf{Surv})}{Z_M}
    =
    \frac{m_M^{(l)}(s)\, \psi_M^{(l)}(s)}{Z_M},
\end{align}
which is~\eqref{ptnb:eq:conditioned-alpha}.

For (iii), substitute~\eqref{ptnb:eq:conditioned-kernel} into the square root to
obtain~\eqref{ptnb:eq:conditioned-geom-kernel}.  Then apply Lemma~\ref{ptnb:lem:bc-chain-rule} to the
conditioned trajectory laws~\eqref{ptnb:eq:conditioned-path-measure} on the ordinary-state graph.  This
gives~\eqref{ptnb:eq:conditioned-bc-kernel} directly.

For (iv), expand the Bhattacharyya coefficient of the conditioned laws:
\begin{align}
    \mathrm{BC}\!\bigl(\Pi_{A,\mathrm{surv}}, \Pi_{B,\mathrm{surv}}\bigr)
    &=
    \sum_{(\zeta_L,\ldots,\zeta_0)\in\mathsf{Surv}}
    \sqrt{\frac{\Pi_A(\zeta_L,\ldots,\zeta_0)}{Z_A}\,
    \frac{\Pi_B(\zeta_L,\ldots,\zeta_0)}{Z_B}} \\
    &=
    \frac{1}{\sqrt{Z_A Z_B}}
    \sum_{(\zeta_L,\ldots,\zeta_0)\in\mathsf{Surv}}
    \sqrt{\Pi_A(\zeta_L,\ldots,\zeta_0)\, \Pi_B(\zeta_L,\ldots,\zeta_0)}.
\end{align}
Grouping the last sum by the ordinary terminal state $\zeta_0=s$ and using the definition of
$g^{(0)}(s)$ exactly as in the proof of Theorem~\ref{ptnb:thm:per-pixel} yields~\eqref{ptnb:eq:ordinary-bc},
hence~\eqref{ptnb:eq:conditioned-bc-posthoc}.  Equation~\eqref{ptnb:eq:conditioned-h2} is then the identity
$\mathrm{H}^2 = 1 - \mathrm{BC}$ applied to the conditioned trajectory laws.
\end{proof}

\paragraph{Implementation consequence.}
Theorem~\ref{ptnb:thm:conditioned-survival} gives two exactly equivalent computations of the conditioned
scalar:
\begin{enumerate}
    \item[(a)] \emph{Kernel-side conditioning:} replace each row by the reweighted
        kernel~\eqref{ptnb:eq:conditioned-kernel}, form the geometric mean~\eqref{ptnb:eq:conditioned-geom-kernel},
        and run the ordinary Hellinger backward pass on $\mathcal{S}^{(L)} \times \cdots \times
            \mathcal{S}^{(0)}$.
    \item[(b)] \emph{Post-hoc normalization:} run the original pass once, extract the ordinary overlap
        $\mathrm{BC}_{\mathrm{ord}}^{(0)} = \sum_{s \in \mathcal{S}^{(0)}} g^{(0)}(s)$ and the survival
        masses $Z_A, Z_B$, and apply~\eqref{ptnb:eq:conditioned-h2}.
\end{enumerate}
The first route makes the conditioned trajectory law explicit; the second is cheaper when only the
scalar is needed, but for deep networks it becomes numerically infeasible when the survival
probability $Z_M$ of either model $M \in \{A,B\}$ gets too small.

\paragraph{Conditioned disagreement map.}
Applying Theorem~\ref{ptnb:thm:per-pixel} to the conditioned kernels $T_{M,\mathrm{surv}}^{(l)}$ yields
the conditioned disagreement map
\begin{align}\label{ptnb:eq:conditioned-per-pixel-h2}
    \mathrm{H}_{\mathrm{surv}}^2(s)
    \;\coloneqq\;
    \tfrac{1}{2}\bigl(m_{A,\mathrm{surv}}^{(0)}(s) + m_{B,\mathrm{surv}}^{(0)}(s)\bigr)
    -
    g_{\mathrm{surv}}^{(0)}(s),
    \qquad s \in \mathcal{S}^{(0)},
\end{align}
which is non-negative and satisfies $\sum_{s \in \mathcal{S}^{(0)}} \mathrm{H}_{\mathrm{surv}}^2(s)
    = \mathrm{H}_{\mathrm{surv}}^2$.

\section{Randomization Bounds}

\begin{lemmaenv}[Alive-set mismatch, \LeanResult{ptnb:lem:alive-mismatch}]\label{ptnb:lem:alive-mismatch}
Suppose a dead state sends all its mass to the cemetery. Every state in the
alive symmetric difference then has zero geometric-mean mass to ordinary
successors:
$\sum_{s'\in\mathcal S^{(l-1)}}\widetilde T^{(l)}(s_i,s')=0$.
If alive states have no cemetery leakage, as for the stated nondegenerate
Stopping-Game kernel, its full geometric-mean row is also zero.\end{lemmaenv}

\begin{proof}
If $i$ is alive in $A$ and dead in $B$, then $T_B^{(l)}(s_i,s')=0$ for every
ordinary successor, so each ordinary geometric mean is zero. The same argument
applies with $A$ and $B$ exchanged. If the alive row also has zero cemetery mass,
the cemetery geometric mean vanishes. With alive-state leakage this last
conclusion does not follow; the ordinary-state conclusion still holds.
\end{proof}

Fix a ReLU MLP with hidden widths $n_1,\ldots,n_L$, an input $x$, and start state $s_0$ at the
target output.  Let $B$ replace $W^{(l^*)}$ by an i.i.d.\ matrix from a continuous distribution
symmetric about zero, and assume that for each $l \ge l^*$ and every neuron $i$,
$\mathbb{P}_B\bigl(i \in \mathrm{alive}^{(l)}_B \mid \text{kernels above layer } l\bigr) \le
    1-\delta$ for some $\delta \in (0,\tfrac12]$.  Then $T_A^{(l)} = T_B^{(l)}$ for $l < l^*$; at $l =
    l^*$ the row weights differ and the alive sets may differ; for $l > l^*$ the perturbation
propagates through the activations.  The same reduction applies to convolutional and transformer
blocks interpreted as single \ac{ptnbLMP} layers.  Write $\eta \coloneqq 1 - \delta \in
    \bigl[\tfrac{1}{2}, 1\bigr)$ and $N \coloneqq L - l^* + 1$.

\begin{theoremenv}[Randomization in the Stopping Game, \LeanResult{ptnb:thm:rand-adf}]\label{ptnb:thm:rand-adf}\label{thm:rand_stopping}
$\mathbb{E}_B\bigl[\mathrm{BC}^{(0)}_{\mathrm{ord}}(\Pi_A, \Pi_B)\bigr] \leq \eta^{N}$.  Hence, for
every $z > 0$, the conditioned-survival distance of Theorem~\ref{ptnb:thm:conditioned-survival}
satisfies $\mathbb{E}_B\bigl[\mathrm{H}^2_{\mathrm{surv}}\,\1\{Z_A Z_B \allowbreak \geq z\}\bigr] \geq
    \mathbb{P}_B(Z_A Z_B \geq z) - \eta^{N}/\sqrt{z}$, while the plain distance is $\mathrm{H}^2 = 1 -
    \mathrm{BC}^{(0)}_{\mathrm{ord}} - g^{(0)}(\dagger)$.\end{theoremenv}

\begin{proof}
At a layer $l \geq l^*$, ordinary overlap mass at a state dead in both models moves into the common
cemetery, mass on the alive symmetric difference vanishes (Lemma~\ref{ptnb:lem:alive-mismatch}), and on
states alive in both models the row mass is at most $1$ by Cauchy--Schwarz.  Hence
$\mathrm{BC}^{(l-1)}_{\mathrm{ord}} \leq \sum_{s \in \mathcal{S}^{(l)}} g^{(l)}(s)\,\1\{s \in
    \mathrm{alive}^{(l)}_B\}$, and, since $g^{(l)}$ is fixed by the kernels above layer~$l$, the
conditional bound gives $\mathbb{E}_B[\mathrm{BC}^{(l-1)}_{\mathrm{ord}}] \leq
    \eta\,\mathbb{E}_B[\mathrm{BC}^{(l)}_{\mathrm{ord}}]$.  Iterating from
$\mathrm{BC}^{(L)}_{\mathrm{ord}} = 1$ over $l = L, \ldots, l^*$ and using that
$\mathrm{BC}^{(l)}_{\mathrm{ord}}$ is non-increasing below $l^*$ yields the bound.  The
$\mathrm{H}_{\mathrm{surv}}$ statement follows from $\mathrm{H}^2_{\mathrm{surv}} = 1 -
    \mathrm{BC}^{(0)}_{\mathrm{ord}}/\sqrt{Z_A Z_B} \geq 1 - \mathrm{BC}^{(0)}_{\mathrm{ord}}/\sqrt{z}$
on $\{Z_A Z_B \geq z\}$ (Theorem~\ref{ptnb:thm:conditioned-survival}), the plain one from
$\mathrm{BC}^{(0)} = \mathrm{BC}^{(0)}_{\mathrm{ord}} + g^{(0)}(\dagger)$.
\end{proof}

\begin{remark}[Cemetery floor]\label{ptnb:rem:adf-cemetery-floor}
Theorem~\ref{ptnb:thm:rand-adf} bounds only the ordinary-state overlap.  Overlap mass that reaches a
neuron dead in both models enters the cemetery in both,
\[
g^{(l-1)}(\dagger) = g^{(l)}(\dagger) +
    \sum_{s \in \mathrm{dead}^{(l)}_A \cap \mathrm{dead}^{(l)}_B} g^{(l)}(s),
\]
and is never thinned
again, so plain $\mathrm{H}$ saturates below one.  In a mean-field approximation with uniform dead
fraction $q$ under $A$ and alive probability $p$ under $B$, the ordinary overlap shrinks by
$(1-q)p$ per layer while a fraction $q(1-p)$ is absorbed, so $g^{(0)}(\dagger) \to
    q(1-p)/(1-(1-q)p)$ as $N \to \infty$; for illustrative values $q = 0.2$, $p = \tfrac{1}{2}$ this
gives $\mathrm{H} \approx 0.91$, the same order as the observed $\drFloorMin$--$\drFloorMax$
(Table~\ref{ptnb:tab:double-random-hellinger}).  The conditioned-survival distance removes the cemetery
floor exactly, and $\mathbb{E}_B[\mathrm{H}^2_{\mathrm{surv}} \mid Z_A Z_B \geq z] \allowbreak\to 1$ for every
fixed $z > 0$ with $\mathbb{P}_B(Z_A Z_B \geq z)$ bounded away from zero.  For the Routing Game the
activation-weighted kernel concentrates $g^{(l)}$ on neurons with largest $|a W|$, which rarely
coincide with $\mathrm{dead}_A \cap \mathrm{dead}_B$, so $\mathrm{H}$ is nearly identical to
$\mathrm{H}_{\mathrm{surv}}$ in practice.
\end{remark}

\begin{theoremenv}[Randomization in the Routing Game, \LeanResult{ptnb:thm:rand-routing}]\label{ptnb:thm:rand-routing}\label{thm:rand_routing}
Under the same setup without the alive-set assumption, take the \ac{rg} kernel at temperature $\tau
    > 0$ and assume that for every $l \geq l^*$ each state alive in both models has row mass
$\sum_{s'}\widetilde{T}^{(l)}(s, s') \leq \bar{b}$ for some $0 \leq \bar{b} < 1$.  Then
$\mathrm{BC}^{(0)}_{\mathrm{ord}}(\Pi_A, \Pi_B) \leq \bar{b}^{N}$.\end{theoremenv}

\begin{proof}
Since $a^{(l-1)}_{M,i} = 0$ for $i \notin \mathrm{alive}^{(l-1)}_M$, the Gibbs
policy~\eqref{ptnb:eq:rg-T} sends no mass to dead neurons, so for $1 \le l < L$ the overlap $g^{(l)}$
vanishes outside $\mathrm{alive}^{(l)}_A \cap \mathrm{alive}^{(l)}_B$, and the start state is alive
in both, in the symmetric difference or dead in both. In the latter two cases its
ordinary overlap contribution is zero by the ordinary-state part of Lemma~\ref{ptnb:lem:alive-mismatch}. Hence
$\mathrm{BC}^{(l-1)}_{\mathrm{ord}} \leq \bar{b}\,\mathrm{BC}^{(l)}_{\mathrm{ord}}$ for every $l
    \geq l^*$, with no alive-set factor (the alive mismatch at layer $l-1$ is already part of the row
masses at layer $l$); iterate as in Theorem~\ref{ptnb:thm:rand-adf}.  The assumption reflects the
structure of the \ac{rg}: each act-to-act row factors as the source-deadness indicator
$\1\{z_j^{(l)} > 0\}$, the constant $\tfrac12$ uniform $\kappa$-branch, and the Gibbs policy
$\pi_{M,\kappa}^{(l)}(\,\cdot\mid j) \propto [a^{(l-1)}_M[\,\cdot\,]\,W^{(l,\kappa)}_M[\,\cdot,
    j]]^{1/\tau}$.  At $l=l^*$ the two policies share $a^{(l^*-1)}$, while their weights may differ; above that layer their activation factors may also differ. These changes can reduce overlap, but a continuous randomization law alone does not imply distinct normalized rows almost surely. For example, if each sign branch has just one positive-mass ordinary route, changing its weight without changing its sign leaves that branch's normalized policy unchanged. The strict conditional row-overlap bound $\bar b<1$ is therefore an explicit hypothesis of this theorem; its applicability requires a separate argument for the architecture and randomization law, or an empirical assessment. The equality case of Cauchy--Schwarz explains strict contraction when the completed normalized rows actually differ.
\end{proof}

